\documentclass[12pt,a4paper]{article}

\usepackage[T2A]{fontenc}
\usepackage[utf8]{inputenc}
\usepackage[russian]{babel}

\usepackage{amsmath,amssymb,mathtools}
\usepackage{geometry}
\usepackage{enumitem}
\usepackage{booktabs}
\usepackage{longtable}
\usepackage{tabularx}
\usepackage{xcolor}
\usepackage{hyperref}
\usepackage{microtype}
\usepackage{array}
\usepackage{fancyhdr}
\usepackage{setspace}
\usepackage{titlesec}
\usepackage{multicol}
\usepackage{pdflscape}

\geometry{margin=2cm}
\setstretch{1.08}
\setlist[itemize]{leftmargin=1.5em}
\setlist[enumerate]{leftmargin=1.7em}

\definecolor{accent}{RGB}{42,84,140}
\definecolor{darkaccent}{RGB}{26,52,92}
\definecolor{textgray}{RGB}{70,70,70}
\definecolor{lightline}{RGB}{210,210,210}

\hypersetup{
    colorlinks=true,
    linkcolor=darkaccent,
    urlcolor=blue!50!black,
    citecolor=darkaccent
}

\pagestyle{fancy}
\fancyhf{}
\fancyhead[L]{\textcolor{textgray}{Водный баланс, дегидратация и ионы у собак и кошек}}
\fancyhead[R]{\textcolor{textgray}{Руководство для ветврача}}
\fancyfoot[C]{\thepage}

\titleformat{\section}{\large\bfseries\color{darkaccent}}{\thesection.}{0.6em}{}
\titleformat{\subsection}{\normalsize\bfseries\color{accent}}{\thesubsection.}{0.6em}{}
\titleformat{\subsubsection}{\normalsize\bfseries}{\thesubsubsection.}{0.5em}{}

\newcommand{\important}[1]{\par\medskip\noindent\textbf{\color{darkaccent}Ключевая мысль.} #1\par\medskip}
\newcommand{\term}[1]{\textbf{#1}}
\newcommand{\na}{\ensuremath{Na^+}}
\newcommand{\ka}{\ensuremath{K^+}}
\newcommand{\cl}{\ensuremath{Cl^-}}
\newcommand{\ca}{\ensuremath{Ca^{2+}}}
\newcommand{\mg}{\ensuremath{Mg^{2+}}}
\newcommand{\phos}{\ensuremath{P}}
\newcommand{\hco}{\ensuremath{HCO_3^-}}
\newcommand{\ica}{\ensuremath{iCa}}
\newcommand{\img}{\ensuremath{iMg}}
\newcommand{\ag}{\ensuremath{AG}}
\newcommand{\sid}{\ensuremath{SID}}

\begin{document}

\begin{center}
    {\LARGE \textbf{Практическое руководство для ветеринарного врача}}\\[0.35em]
    {\Large \textbf{Водный баланс, типы дегидратации и нарушения основных ионов}}\\[0.35em]
    {\large Собаки и кошки}\\[0.7em]
    \textcolor{textgray}{Современный клинический конспект для приёма, стационара и ОРИТ}
\end{center}

\vspace{0.6em}
\noindent\rule{\textwidth}{0.5pt}

\pdfbookmark[1]{Оглавление}{toc}
\tableofcontents
\newpage

\section{Зачем это руководство}
Нарушения водного баланса и электролитов у собак и кошек редко существуют изолированно. На практике врач почти всегда имеет дело с комбинацией:
\begin{itemize}
    \item гипоперфузии или гиповолемии;
    \item общей дегидратации или, наоборот, перегрузки жидкостью;
    \item сдвигов натрия, калия, хлора, кальция, магния и фосфора;
    \item кислотно-основных нарушений;
    \item первичного заболевания, которое эти сдвиги вызывает.
\end{itemize}

\important{У пациента нужно разделять три разные задачи: восстановить перфузию, оценить общий водный баланс и отдельно скорректировать опасные электролитные нарушения.}

\section{Базовая патофизиология водного баланса}
Общая вода организма распределяется между внутриклеточным и внеклеточным секторами. Для клиники особенно важны:
\begin{itemize}
    \item \term{внутрисосудистый объём} --- определяет перфузию и артериальное давление;
    \item \term{интерстициальная жидкость} --- отражает общую гидратацию;
    \item \term{внутриклеточная вода} --- чувствительна к изменениям осмоляльности, прежде всего к натрию.
\end{itemize}

Натрий --- основной катион внеклеточной жидкости и главный детерминант её тоничности. Калий --- основной внутриклеточный катион. Хлор тесно связан с кислотно-основным состоянием. Ионизированный кальций и магний определяют нервно-мышечную возбудимость и стабильность проводящей системы сердца. Фосфор особенно важен для энергетического обмена, мембран, эритроцитов и нервно-мышечной функции.

\subsection{Что влияет на водный баланс}
\textbf{Со стороны внешней среды:}
\begin{itemize}
    \item температура;
    \item влажность;
    \item доступ к воде;
    \item физическая нагрузка;
    \item солевая нагрузка.
\end{itemize}

\textbf{Со стороны организма:}
\begin{itemize}
    \item анорексия, снижение потребления воды;
    \item рвота, диарея, секвестрация жидкости в ЖКТ;
    \item полиурия, осмотический диурез, почечные потери;
    \item лихорадка, тахипноэ, ожоги;
    \item потери в третье пространство;
    \item эндокринные причины: сахарный диабет, гипоадренокортицизм, несахарный диабет;
    \item обструкция мочевыводящих путей и нарушения выведения жидкости.
\end{itemize}

\section{Первичная клиническая логика}
\subsection{Первый вопрос: есть ли гипоперфузия?}
При поступлении пациента сначала отвечают на вопрос, есть ли гемодинамически значимая гиповолемия или шок.

\textbf{Ранние признаки гипоперфузии:}
\begin{itemize}
    \item тахикардия;
    \item удлинение CRT;
    \item ухудшение качества пульса;
    \item жажда;
    \item беспокойство или депрессия;
    \item снижение диуреза;
    \item холодные конечности.
\end{itemize}

\textbf{Поздние признаки:}
\begin{itemize}
    \item слабый пульс;
    \item аритмии;
    \item выраженная слабость;
    \item нарушение сознания;
    \item гипотермия;
    \item коллапс;
    \item цианоз.
\end{itemize}

\important{Если пациент гиповолемичен, сначала восстанавливают перфузию изотоническими кристаллоидами, а уже потом решают вопрос медленной коррекции натрия и свободной воды.}

\subsection{Что спросить в анамнезе}
\begin{itemize}
    \item аппетит и жажда;
    \item рвота, диарея, регургитация;
    \item полиурия, полидипсия, олигурия, анурия;
    \item странгурия, обструкция, травма мочевых путей;
    \item лихорадка, тепловая нагрузка, тахипноэ;
    \item наличие отёков, выпотов, асцита;
    \item предшествующая инфузионная терапия;
    \item диуретики, инсулин, стероиды;
    \item болезни почек, сердца, печени, надпочечников;
    \item диабет, панкреатит, сепсис.
\end{itemize}

\subsection{Что взять в первые 15--30 минут}
\begin{itemize}
    \item PCV/TS или альбумин/общий белок;
    \item электролиты: \na, \ka, \cl, по возможности \ica\ и \img;
    \item глюкоза;
    \item креатинин, мочевина;
    \item газовый состав крови / КЩС;
    \item лактат по возможности;
    \item общий анализ мочи и оценка удельного веса мочи;
    \item ЭКГ при подозрении на значимые сдвиги \ka\ или \ca.
\end{itemize}

\section{Мониторинг}
\subsection{Клинический минимум}
\begin{enumerate}
    \item осмотр 2--3 раза в сутки, у тяжёлых пациентов чаще;
    \item масса тела минимум 1 раз в сутки, у критических нередко 2 раза в сутки;
    \item учёт входов и выходов;
    \item оценка диуреза;
    \item электролиты по показаниям и при активной коррекции;
    \item глюкоза, белки/альбумин, КЩС;
    \item ЭКГ при подозрении на значимые сдвиги \ka\ или \ca.
\end{enumerate}

\subsection{Диурез}
Практический ориентир для стационарного пациента --- около \(1\ \text{мл/кг/ч}\) и выше, но интерпретировать его нужно совместно с:
\begin{itemize}
    \item артериальным давлением;
    \item объёмом инфузии;
    \item креатинином и мочевиной;
    \item гидратацией;
    \item вероятностью обструкции.
\end{itemize}

\subsection{Почему важен вес}
Изменение массы тела у госпитализированного животного обычно отражает изменение воды, а не ткани. Быстрое увеличение веса на фоне инфузии --- тревожный сигнал возможной перегрузки жидкостью или накопления в третьем пространстве.

\subsection{Шаблон динамического мониторинга}
\begin{longtable}{>{\raggedright\arraybackslash}p{0.28\textwidth} >{\raggedright\arraybackslash}p{0.22\textwidth} >{\raggedright\arraybackslash}p{0.21\textwidth} >{\raggedright\arraybackslash}p{0.21\textwidth}}
\toprule
\textbf{Параметр} & \textbf{Как часто} & \textbf{Когда особенно нужен} & \textbf{Тревожные признаки} \\
\midrule
Осмотр, ментальный статус, CRT, пульс & каждые 1--6 ч & шок, нестабильный пациент, активная коррекция & нарастание тахикардии, слабый пульс, ухудшение сознания \\
Масса тела & 1--2 раза/сут & стационар, риск перегрузки, ОПП/ХБП & быстрый прирост массы \\
Диурез & почасово у тяжёлых & ОРИТ, ОПП, обструкция, гиперкалиемия & олигурия/анурия, отсутствие ответа на инфузию \\
\na, \ka, \cl & каждые 2--12 ч & активная коррекция, ДКА, Аддисон, ОПП, обструкция & быстрый сдвиг Na, высокий K, тяжёлая гипохлоремия \\
\ica\ / \img & по показаниям, нередко каждые 4--12 ч & тремор, аритмии, панкреатит, сепсис, переливания & симптомная гипокальциемия, рефрактерная гипокалиемия \\
Глюкоза & каждые 1--6 ч & инсулин, ДКА, гиперкалиемия, критические пациенты & гипо-/гипергликемия \\
КЩС / газы крови & каждые 4--12 ч & тяжёлый ацидоз/алкалоз, ДКА, сепсис & нарастание ацидоза, гиперхлоремический ацидоз \\
ЭКГ & непрерывно или серийно & гиперкалиемия, симптомная гипокальциемия & AV-блокада, широкий QRS, тяжёлая брадикардия \\
\bottomrule
\end{longtable}

\section{Как считать водный баланс}
\[
\text{Водный баланс} = \text{входы} - \text{выходы}
\]

\textbf{Входы:}
\begin{itemize}
    \item инфузия;
    \item выпитая вода;
    \item влажный корм;
    \item лекарственные растворы;
    \item болюсы, промывки, шприцевые насосы, энтеральное питание.
\end{itemize}

\textbf{Выходы:}
\begin{itemize}
    \item моча;
    \item рвота;
    \item диарея;
    \item аспираты и дренажи;
    \item крупные выпоты при удалении;
    \item оценочные потери при лихорадке и тахипноэ.
\end{itemize}

\textbf{Интерпретация:}
\begin{itemize}
    \item положительный баланс: введено больше, чем выведено;
    \item нулевой баланс: входы и выходы сопоставимы;
    \item отрицательный баланс: потери превышают поступление.
\end{itemize}

\important{Положительный водный баланс без улучшения перфузии и без адекватного диуреза --- повод немедленно пересмотреть инфузионный план.}

\section{Типы дегидратации}
\subsection{Изотоническая дегидратация}
Потери воды и натрия близки по составу к внеклеточной жидкости. Страдает прежде всего внеклеточный сектор и объём циркулирующей крови.

\textbf{Типичные причины:}
\begin{itemize}
    \item рвота;
    \item диарея;
    \item кровопотеря;
    \item секвестрация в ЖКТ;
    \item потери в третье пространство.
\end{itemize}

\subsection{Гипертоническая дегидратация}
Потеря свободной воды превышает потерю натрия. Внеклеточная жидкость становится гиперосмолярной, вода выходит из клеток.

\textbf{Типичные причины:}
\begin{itemize}
    \item ограничение доступа к воде;
    \item несахарный диабет;
    \item осмотический диурез;
    \item лихорадка;
    \item тахипноэ;
    \item тепловой стресс.
\end{itemize}

\textbf{Клинический риск:} неврологические нарушения из-за клеточной дегидратации мозга.

\subsection{Гипотоническая дегидратация}
Потеря натрия превышает потерю воды. Осмоляльность внеклеточной жидкости снижается, вода переходит в клетки.

\textbf{Типичные причины:}
\begin{itemize}
    \item гипоадренокортицизм;
    \item тяжёлые ЖКТ-потери;
    \item некоторые нефропатии;
    \item избыточное введение гипотонических растворов.
\end{itemize}

\textbf{Клинический риск:} сочетание гиповолемии с риском клеточного отёка, включая отёк мозга.

\section{Натрий}
\subsection{Что означает концентрация натрия}
Концентрация \na\ --- это прежде всего показатель соотношения натрия к воде. Поэтому гипер- и гипонатриемия почти всегда отражают нарушение водного баланса, а не просто изменение общего количества натрия в организме.

\subsection{Регуляция}
Основные регуляторы:
\begin{itemize}
    \item жажда;
    \item АДГ;
    \item РААС;
    \item симпатическая нервная система.
\end{itemize}

При снижении эффективного артериального объёма активируются ренин, ангиотензин II и альдостерон; растёт реабсорбция натрия, а АДГ повышает удержание воды.

\subsection{Гипернатриемия}
\textbf{Основные варианты по объёмному статусу:}
\begin{itemize}
    \item \textbf{гиповолемическая:} потери воды при ЖКТ-потерях, почечных потерях, лихорадке, тахипноэ;
    \item \textbf{эуволемическая:} потеря свободной воды, несахарный диабет;
    \item \textbf{гиперволемическая:} гипертонические натрийсодержащие растворы, солевые интоксикации.
\end{itemize}

\textbf{Клиника:}
\begin{itemize}
    \item вялость;
    \item слабость;
    \item атаксия;
    \item тремор;
    \item судороги;
    \item кома.
\end{itemize}

\subsection{Гипонатриемия}
\textbf{Основные варианты по объёмному статусу:}
\begin{itemize}
    \item \textbf{гиповолемическая:} потеря \na\ при рвоте, диарее, диуретиках, гипоадренокортицизме;
    \item \textbf{эуволемическая:} SIADH-подобные состояния, гипотонические инфузии;
    \item \textbf{гиперволемическая:} водная задержка при сердечной и почечной недостаточности, иногда при тяжёлой болезни печени.
\end{itemize}

\textbf{Клиника:}
\begin{itemize}
    \item слабость;
    \item депрессия;
    \item тошнота/рвота;
    \item атаксия;
    \item тремор;
    \item судороги;
    \item кома.
\end{itemize}

\subsection{Главный принцип коррекции натрия}
\important{Быстрая коррекция диснатриемии может быть опаснее самого нарушения.}

\textbf{Практические правила:}
\begin{itemize}
    \item сначала стабилизируют гемодинамику;
    \item затем корректируют водный компонент и натрий;
    \item обычно стараются не изменять \na\ быстрее, чем примерно на \(0{,}5\ \text{ммоль/л/ч}\);
    \item при хронических нарушениях ориентируются также на суточный предел порядка \(8\text{--}12\ \text{ммоль/л/сут}\), если нет немедленной неврологической угрозы.
\end{itemize}

\subsection{Дефицит свободной воды}
После устранения гиповолемии дефицит свободной воды можно оценить:
\[
\text{ДСВ} = 0{,}6 \cdot m(\text{кг}) \cdot \left(\frac{Na_{\text{изм}}}{Na_{\text{норма}}} - 1\right)
\]

Это расчётная отправная точка, а не абсолютная истина. Реальная коррекция всегда опирается на серийный контроль натрия и клинический ответ.

\subsection{Гипонатриемия и гипергликемия}
При выраженной гипергликемии натрий может казаться ниже из-за водного перераспределения. Практически полезно считать скорректированный натрий:
\[
Na_{\text{corr}} \approx Na_{\text{изм}} + 1.6 \cdot \frac{(\text{глюкоза мг/дл} - 100)}{100}
\]
Если глюкоза выражена в ммоль/л, формулу нужно адаптировать под единицы.

\section{Калий}
\subsection{Клиническое значение}
Калий --- электролит немедленной опасности, потому что быстро меняет мембранный потенциал и проводимость миокарда.

\subsection{Гиперкалиемия}
\textbf{Основные причины:}
\begin{itemize}
    \item снижение почечного выведения;
    \item острая или хроническая почечная недостаточность;
    \item уретральная обструкция;
    \item разрыв мочевого пузыря;
    \item гипоадренокортицизм;
    \item ацидоз;
    \item инсулиновая недостаточность.
\end{itemize}

\textbf{Типичные ЭКГ-изменения:}
\begin{itemize}
    \item высокий пикообразный зубец \(T\);
    \item уплощение или исчезновение зубца \(P\);
    \item удлинение \(PR\);
    \item расширение \(QRS\);
    \item AV-блокада;
    \item синусоидальная форма при тяжёлых случаях.
\end{itemize}

\subsection{Экстренное лечение гиперкалиемии}
\begin{enumerate}
    \item \textbf{Стабилизировать миокард:} кальция глюконат внутривенно медленно под ЭКГ-контролем.
    \item \textbf{Сместить \ka\ в клетку:} инсулин короткого действия с декстрозой; при выраженном метаболическом ацидозе иногда дополнительно бикарбонат.
    \item \textbf{Удалить \ka\ из организма и устранить причину:}
    \begin{itemize}
        \item инфузия и восстановление перфузии;
        \item восстановление диуреза;
        \item устранение обструкции;
        \item диализ при рефрактерных случаях.
    \end{itemize}
\end{enumerate}

\important{Кальций улучшает электрическую стабильность миокарда, но не снижает общий калий в организме.}

\subsection{Гипокалиемия}
\textbf{Причины:}
\begin{itemize}
    \item недостаточное поступление;
    \item рвота, диарея;
    \item почечные потери;
    \item диуретики;
    \item алкалоз;
    \item осмотический диурез;
    \item гиперальдостеронизм.
\end{itemize}

\textbf{Клиника:}
\begin{itemize}
    \item слабость;
    \item анорексия;
    \item парез;
    \item ухудшение нервно-мышечной проводимости;
    \item у кошек --- вентрофлексия шеи.
\end{itemize}

\subsection{Коррекция гипокалиемии}
Для внутривенного введения \ka\ обычно придерживаются верхнего предела порядка:
\[
0{,}5\ \text{ммоль/кг/ч}
\]

Стабильным пациентам как можно раньше подключают пероральные препараты калия.

\important{Если гипокалиемия корректируется плохо, нужно проверить магний: дефицит \mg\ часто делает гипокалиемию рефрактерной.}

\section{Хлор}
\subsection{Почему хлор важен}
Хлор тесно связан с кислотно-основным состоянием и не должен рассматриваться как второстепенный электролит.

\subsection{Гиперхлоремия}
\textbf{Частые ситуации:}
\begin{itemize}
    \item гипернатриемия;
    \item большие объёмы \(0{,}9\%\) NaCl;
    \item гиперхлоремический метаболический ацидоз.
\end{itemize}

\subsection{Гипохлоремия}
\textbf{Частые ситуации:}
\begin{itemize}
    \item рвота и потери HCl;
    \item обструкция верхних отделов ЖКТ;
    \item гипохлоремический метаболический алкалоз.
\end{itemize}

\important{У пациента с рвотой, высоким \hco\ и низким \cl\ надо думать не просто о дегидратации, а о chloride-responsive alkalosis.}

\section{Кальций}
\subsection{Фракции кальция}
\begin{itemize}
    \item ионизированный --- около \(55\%\);
    \item связанный с анионами --- около \(10\%\);
    \item связанный с белками --- около \(35\%\).
\end{itemize}

Биологически активен именно ионизированный кальций. У критических пациентов ориентироваться лучше на \ica, а не только на общий кальций.

\subsection{Гиперкальциемия}
\textbf{Важные причины:}
\begin{itemize}
    \item гиперкальциемия злокачественных новообразований;
    \item первичный гиперпаратиреоз;
    \item аденокарцинома апокриновых желёз анальных мешков;
    \item множественная миелома;
    \item ХБП;
    \item гипоадренокортицизм;
    \item гранулематозные болезни;
    \item интоксикация витамином D;
    \item идиопатическая гиперкальциемия, особенно у кошек.
\end{itemize}

\textbf{Клиника:}
\begin{itemize}
    \item анорексия;
    \item вялость;
    \item полиурия/полидипсия;
    \item слабость;
    \item рвота;
    \item запор;
    \item снижение концентрационной способности почек.
\end{itemize}

\textbf{Лечение:}
\begin{itemize}
    \item регидратация;
    \item после регидратации --- фуросемид;
    \item при необходимости --- бисфосфонаты;
    \item при тяжёлых рефрактерных случаях --- диализ;
    \item глюкокортикоиды --- желательно после получения диагностического материала, если подозревается неоплазия.
\end{itemize}

\subsection{Гипокальциемия}
\textbf{Причины:}
\begin{itemize}
    \item гипопаратиреоз;
    \item эклампсия;
    \item панкреатит;
    \item сепсис;
    \item этиленгликоль;
    \item цитратная нагрузка при массивных трансфузиях;
    \item ДВС-синдром;
    \item гиповитаминоз D;
    \item тяжёлое критическое состояние.
\end{itemize}

\textbf{Клиника:}
\begin{itemize}
    \item нервно-мышечная возбудимость;
    \item тремор;
    \item фасцикуляции;
    \item поведенческие изменения;
    \item слабость;
    \item тетания;
    \item судороги;
    \item иногда стридор и дыхательные нарушения.
\end{itemize}

\textbf{Лечение:}
\begin{itemize}
    \item кальция глюконат \(10\%\) внутривенно медленно под ЭКГ-контролем;
    \item затем --- поддерживающая терапия и лечение первичной причины;
    \item в дальнейшем ориентироваться на \ica.
\end{itemize}

\section{Магний}
\subsection{Клиническое значение}
Магний влияет на:
\begin{itemize}
    \item нервно-мышечную передачу;
    \item возбудимость миокарда;
    \item внутриклеточный транспорт калия;
    \item взаимодействие с кальцием.
\end{itemize}

\subsection{Когда подозревать гипомагниемию}
\begin{itemize}
    \item рефрактерная гипокалиемия;
    \item аритмии;
    \item тремор;
    \item повышенная нервно-мышечная возбудимость;
    \item длительная полиурия;
    \item ДКА;
    \item сепсис;
    \item длительные ЖКТ-потери.
\end{itemize}

\textbf{Клиника гипомагниемии:}
\begin{itemize}
    \item анорексия;
    \item тремор;
    \item слабость;
    \item аритмии.
\end{itemize}

\textbf{Клиника гипермагниемии:}
\begin{itemize}
    \item слабость;
    \item гипорефлексия;
    \item угнетение дыхания;
    \item брадикардия.
\end{itemize}

\section{Фосфор}
\subsection{Почему стоит добавить фосфор в пособие}
Фосфор критичен для синтеза АТФ, работы мембран, эритроцитов, мышц и нервной системы. У многих пациентов с тяжёлой внутренней болезнью проблема не ограничивается Na/K/Cl.

\subsection{Гипофосфатемия}
\textbf{Типичные причины:}
\begin{itemize}
    \item ДКА после начала инсулина;
    \item рефидинг;
    \item сепсис;
    \item тяжёлая анорексия;
    \item алкалоз;
    \item усиленная почечная экскреция.
\end{itemize}

\textbf{Клиника:}
\begin{itemize}
    \item слабость;
    \item гемолиз;
    \item снижение сократимости мышц и дыхательной мускулатуры;
    \item неврологические нарушения;
    \item ухудшение общего восстановления.
\end{itemize}

\subsection{Гиперфосфатемия}
\textbf{Типичные причины:}
\begin{itemize}
    \item ОПП/ХБП;
    \item массивное клеточное повреждение;
    \item тяжёлое нарушение выведения.
\end{itemize}

\textbf{Почему это важно:}
\begin{itemize}
    \item влияет на кальций;
    \item поддерживает вторичные метаболические сдвиги;
    \item при высоком произведении \(Ca \times P\) растёт риск минерализации тканей.
\end{itemize}

\subsection{Практика}
При ДКА, рефидинге, тяжёлой анорексии, сепсисе и почечной недостаточности контроль фосфора должен быть рутинным, а не факультативным.

\section{Кислотно-основное состояние}
КЩС интерпретируют вместе с электролитами.

\textbf{Практические связи:}
\begin{itemize}
    \item ацидоз может способствовать выходу \ka\ во внеклеточное пространство;
    \item рвота ведёт к потере HCl и часто даёт гипохлоремический метаболический алкалоз;
    \item большие объёмы \(0{,}9\%\) NaCl могут способствовать гиперхлоремическому метаболическому ацидозу;
    \item бикарбонат применяют избирательно, а не ``по привычке''.
\end{itemize}

\section{Выбор инфузионного раствора}
\begin{longtable}{>{\raggedright\arraybackslash}p{0.19\textwidth} >{\raggedright\arraybackslash}p{0.23\textwidth} >{\raggedright\arraybackslash}p{0.24\textwidth} >{\raggedright\arraybackslash}p{0.24\textwidth}}
\toprule
\textbf{Раствор} & \textbf{Когда уместен} & \textbf{Когда нежелателен} & \textbf{Особенности} \\
\midrule
Сбалансированный изотонический кристаллоид & гиповолемия, регидратация, большинство стационарных пациентов & тяжёлые специфические ситуации, где нужен иной состав по Na/Cl & содержит буферные анионы; обычно быстрее улучшает КЩС, чем \(0{,}9\%\) NaCl \\
\(0{,}9\%\) NaCl & гипохлоремический метаболический алкалоз, потери HCl, отдельные ситуации с натрием/хлором & риск гиперхлоремического ацидоза, перегрузка хлоридом & высокий Cl, может усиливать гиперхлоремию \\
Растворы с декстрозой & свободная вода после стабилизации, часть планов коррекции диснатриемии, гипогликемия & первичная коррекция гиповолемии & не подходят для ресусцитации \\
Гипотонические растворы & после стабилизации, для избранных пациентов под строгим контролем Na & шок, гиповолемия, неясная диснатриемия & быстрое и неосторожное использование опасно \\
Гипертонический NaCl & отдельные реанимационные сценарии, некоторые неврологические пациенты, ограниченная объёмная ресусцитация & хроническая гипернатриемия, обезвоживание без плана мониторинга & требует чёткого понимания цели и частого контроля Na \\
\bottomrule
\end{longtable}

\important{Для первичной коррекции гиповолемии растворы свободной воды и гипотонические растворы не подходят.}

\section{Численные пороги и уровни срочности}
\subsection{Практическая таблица срочности}
\begin{longtable}{>{\raggedright\arraybackslash}p{0.17\textwidth} >{\raggedright\arraybackslash}p{0.17\textwidth} >{\raggedright\arraybackslash}p{0.18\textwidth} >{\raggedright\arraybackslash}p{0.21\textwidth} >{\raggedright\arraybackslash}p{0.18\textwidth}}
\toprule
\textbf{Показатель} & \textbf{Лёгкое отклонение} & \textbf{Умеренное отклонение} & \textbf{Тяжёлое / потенциально опасное} & \textbf{Что делать} \\
\midrule
\na & небольшое отклонение без неврологии & выраженный сдвиг без тяжёлой неврологии & быстрый сдвиг, судороги, кома & оценить объём, часто пересдавать, корректировать медленно \\
\ka & mild hypo-/hyper-K & есть клиника или сопутствующая аритмия & ЭКГ-изменения, тяжёлая брадикардия, широкий QRS & ЭКГ, срочная коррекция \\
\cl & изолированный сдвиг без КЩС-эффекта & есть алкалоз/ацидоз & тяжёлый сдвиг как часть критического состояния & лечить причину, смотреть КЩС и Na \\
\ica & небольшое снижение/повышение без симптомов & тремор, слабость, ЖКТ-признаки & тетания, аритмии, судороги & контролировать \ica, лечить симптомных \\
\mg & лабораторное отклонение без клиники & аритмии, тремор, рефрактерная hypo-K & выраженные нейромышечные или сердечные проявления & срочная коррекция и повторный контроль \\
\phos & mild hypo-/hyper-P & слабость, риск гемолиза, ДКА/рефидинг & выраженная симптоматика, дыхательная слабость, гемолиз & корректировать и мониторировать в динамике \\
\bottomrule
\end{longtable}

\textbf{Примечание.} В этой таблице намеренно не приведены универсальные числовые границы для всех лабораторий: они должны сверяться с референсами конкретной лаборатории, анализатора и клиническим контекстом.

\section{Формулы и быстрые расчёты}
\subsection{Водный баланс}
\[
\text{Водный баланс} = \text{входы} - \text{выходы}
\]

\subsection{Дефицит свободной воды}
\[
\text{ДСВ} = 0{,}6 \cdot m(\text{кг}) \cdot \left(\frac{Na_{\text{изм}}}{Na_{\text{норма}}} - 1\right)
\]

\subsection{Скорректированный натрий при гипергликемии}
\[
Na_{\text{corr}} \approx Na_{\text{изм}} + 1.6 \cdot \frac{(\text{глюкоза мг/дл} - 100)}{100}
\]

\subsection{Анионная разница}
\[
\ag = (\na + \ka) - (\cl + \hco)
\]
Во многих клиниках допустимо использовать и упрощённый вариант без калия:
\[
\ag = \na - (\cl + \hco)
\]

\subsection{Скорректированный хлор}
Полезный приближённый расчёт:
\[
Cl_{\text{corr}} = Cl_{\text{изм}} \cdot \frac{Na_{\text{норма}}}{Na_{\text{изм}}}
\]
Он помогает понять, действительно ли хлор ``сам по себе'' изменён, или его сдвиг в основном вторичен по отношению к натрию.

\subsection{Упрощённый strong ion difference}
\[
\sid \approx (\na + \ka) - \cl
\]
Это упрощение не заменяет полноценный анализ, но может быть полезно для клинического мышления у сложных пациентов.

\subsection{Безопасная скорость изменения натрия}
\[
\Delta Na_{\max} \approx 0{,}5\ \text{ммоль/л/ч}
\]

\subsection{Практический верхний ориентир скорости внутривенного калия}
\[
K_{\max} \approx 0{,}5\ \text{ммоль/кг/ч}
\]

\section{Перегрузка жидкостью}
\subsection{Почему это нужно выделять отдельно}
Неответ на инфузию не всегда означает ``нужно ещё жидкости''. Иногда это уже перегрузка жидкостью, третье пространство или кардиоренальная проблема.

\subsection{Признаки перегрузки}
\begin{itemize}
    \item быстрый прирост массы тела;
    \item тахипноэ, усиление дыхательных усилий;
    \item серозные выделения из носа;
    \item хрипы;
    \item хемоз;
    \item появление или усиление выпотов/асцита;
    \item ухудшение оксигенации;
    \item положительный баланс без улучшения перфузии;
    \item снижение диуреза или отсутствие ответа на инфузию.
\end{itemize}

\subsection{Группа риска}
\begin{itemize}
    \item кошки;
    \item пациенты с ХБП/ОПП;
    \item сердечные пациенты;
    \item гипоальбуминемия;
    \item сепсис и выраженное третье пространство;
    \item очень мелкие животные.
\end{itemize}

\subsection{Что делать}
\begin{enumerate}
    \item немедленно пересмотреть показания к продолжающейся инфузии;
    \item повторно оценить перфузию, АД, массу тела, диурез;
    \item уменьшить скорость или временно остановить инфузию при подозрении на перегрузку;
    \item рассмотреть кислород, визуализацию, кардиологическую оценку;
    \item диуретики только после клинической переоценки, а не автоматически;
    \item у олигурических пациентов пересмотреть риск необходимости декомпрессии, катетеризации, нефрологического или интенсивного ведения.
\end{enumerate}

\section{POCUS при нарушениях водного баланса}
\subsection{Зачем POCUS нужен в этой теме}
Bedside ultrasound позволяет быстрее понять, с чем именно связан нестабильный пациент:
\begin{itemize}
    \item истинная гиповолемия;
    \item скрытое кровотечение или выпот;
    \item плевральный выпот;
    \item признаки отёка лёгких;
    \item выраженное третье пространство;
    \item кардиогенный компонент ухудшения.
\end{itemize}

\important{POCUS не заменяет полный УЗИ-протокол, но очень полезен как быстрый инструмент принятия решений у неотложного пациента.}

\subsection{Что смотреть}
\textbf{AFAST / abdominal POCUS:}
\begin{itemize}
    \item свободная жидкость в брюшной полости;
    \item выраженность выпота;
    \item крупные изменения мочевого пузыря;
    \item косвенные признаки секвестрации жидкости.
\end{itemize}

\textbf{TFAST / thoracic POCUS:}
\begin{itemize}
    \item плевральный выпот;
    \item B-lines как косвенный признак интерстициально-альвеолярного синдрома;
    \item скольжение плевры;
    \item грубые изменения, подозрительные на дыхательную декомпенсацию.
\end{itemize}

\textbf{CPOCUS / cardiac POCUS:}
\begin{itemize}
    \item наличие перикардиального выпота;
    \item грубая субъективная оценка сократимости;
    \item выраженность дилатации камер;
    \item очевидные признаки, что пациент плохо переносит дальнейшую объёмную нагрузку.
\end{itemize}

\subsection{Когда особенно полезен}
\begin{itemize}
    \item шок или неясная гипоперфузия;
    \item слабый ответ на болюсную терапию;
    \item риск перегрузки жидкостью;
    \item подозрение на выпоты;
    \item кардиологический пациент;
    \item пациент с олигурией и нарастающим положительным балансом.
\end{itemize}

\subsection{Чего POCUS не решает}
\begin{itemize}
    \item не отменяет лабораторную оценку электролитов и КЩС;
    \item не заменяет полноценное ЭХО;
    \item не должен использоваться в отрыве от клинической картины и мониторинга.
\end{itemize}

\section{Цели ресусцитации и конечные точки ответа на жидкость}
\subsection{Главная идея}
Инфузионная терапия должна иметь чёткую цель. Вводить жидкость ``пока не закончится пакет'' --- плохая практика.

\subsection{Практические конечные точки}
После болюса или этапа ресусцитации врач должен оценить:
\begin{itemize}
    \item CRT;
    \item качество и наполнение пульса;
    \item ЧСС;
    \item температуру конечностей;
    \item ментальный статус;
    \item артериальное давление;
    \item диурез;
    \item лактат / динамику перфузионных маркеров по возможности;
    \item дыхательный статус;
    \item риск перегрузки жидкостью.
\end{itemize}

\subsection{Когда болюсы стоит остановить}
\begin{itemize}
    \item достигнут явный клинический ответ;
    \item перестал улучшаться пульс, CRT, ментальный статус;
    \item появилась тахипноэ или ухудшилось дыхание;
    \item появился хемоз, серозные выделения, подозрение на перегрузку;
    \item пациент кардиологический или почечный и дальнейший объём начинает нести больше риска, чем пользы.
\end{itemize}

\important{Цель жидкостной ресусцитации --- не ``дать объём'', а восстановить достаточную перфузию при минимально возможной цене по риску перегрузки.}

\section{Replacement, maintenance, ongoing losses}
\subsection{Три компонента инфузионного плана}
Любой жидкостный план удобно делить на три части:
\begin{enumerate}
    \item \textbf{replacement} --- восполнение дефицита;
    \item \textbf{maintenance} --- поддерживающая потребность;
    \item \textbf{ongoing losses} --- продолжающиеся потери.
\end{enumerate}

\subsection{Replacement}
Это:
\begin{itemize}
    \item дефицит гидратации;
    \item внутрисосудистая гиповолемия;
    \item часть ресусцитации у нестабильного пациента.
\end{itemize}

\subsection{Maintenance}
Это базовая суточная потребность в жидкости у пациента без больших дополнительных потерь. Этот компонент нужно уменьшать или пересматривать:
\begin{itemize}
    \item у олигурических пациентов;
    \item при риске перегрузки;
    \item у кошек;
    \item у кардиологических и некоторых нефрологических пациентов.
\end{itemize}

\subsection{Ongoing losses}
Это продолжающиеся потери, которые нельзя игнорировать:
\begin{itemize}
    \item рвота;
    \item диарея;
    \item дренажи;
    \item аспираты;
    \item полиурия;
    \item выпоты при активном удалении;
    \item потери через дыхание и гипертермию.
\end{itemize}

\subsection{Практическая ошибка}
Одна из самых частых ошибок --- посчитать только стартовый дефицит и забыть, что через 4--6 часов пациент уже живёт в другой точке: объём изменился, потери продолжаются, электролиты сдвинулись.

\important{Инфузионный план должен пересчитываться повторно, а не только назначаться при поступлении.}

\section{Preanalytical и laboratory pitfalls}
\subsection{Почему это важно}
До того как интерпретировать ``опасный'' анализ, нужно убедиться, что проба достойна интерпретации.

\subsection{Частые источники ошибок}
\begin{itemize}
    \item гемолиз;
    \item задержка обработки пробы;
    \item неправильное хранение;
    \item загрязнение образца инфузионным раствором;
    \item забор крови из линии, по которой недавно шла инфузия;
    \item липемия;
    \item неправильный антикоагулянт или проблемы с пробиркой;
    \item длительное стояние цельной крови до центрифугирования.
\end{itemize}

\subsection{Что может исказиться}
\begin{itemize}
    \item \ka\ --- ложное повышение при гемолизе;
    \item глюкоза --- снижение при задержке анализа;
    \item натрий / хлор --- искажение при контаминации раствором;
    \item кальций --- ошибки интерпретации при неподходящем образце или отсрочке анализа;
    \item КЩС --- быстро теряет клиническую ценность при неправильной преаналитике.
\end{itemize}

\subsection{Практическое правило}
Если результат не сходится с клиникой, врач обязан подумать:
\begin{itemize}
    \item это настоящий сдвиг?
    \item это артефакт?
    \item нужно ли срочно пересдать?
\end{itemize}

\important{Критическое мышление начинается раньше интерпретации цифры --- с оценки качества пробы.}

\section{Кардиологический пациент и инфузионная терапия}
\subsection{Почему нужен отдельный раздел}
Пациент с сердечной недостаточностью или подозрением на неё может одновременно иметь:
\begin{itemize}
    \item низкий эффективный артериальный объём;
    \item плевральный выпот или асцит;
    \item риск отёка лёгких;
    \item плохую переносимость даже умеренной объёмной нагрузки.
\end{itemize}

\subsection{Что помнить}
\begin{itemize}
    \item выпоты и отёки не означают, что внутрисосудистый объём нормален;
    \item но и не дают права ``лить по стандарту'';
    \item кардиологический пациент требует меньших шагов инфузионной оценки;
    \item болюсы должны быть осторожнее и с очень частым пересмотром;
    \item POCUS, дыхательный мониторинг и масса тела особенно полезны;
    \item ухудшение дыхания на инфузии --- повод немедленно остановиться и пересмотреть план.
\end{itemize}

\subsection{Когда думать о сердечной причине}
\begin{itemize}
    \item тахипноэ непропорционально клинике дегидратации;
    \item ортопноэ;
    \item хрипы или приглушённые тоны;
    \item выпоты;
    \item отсутствие ожидаемого ответа на жидкость;
    \item ухудшение после небольших объёмов инфузии.
\end{itemize}

\section{Печёночный пациент}
\subsection{Что делает этих пациентов особенными}
У печёночного пациента могут одновременно сосуществовать:
\begin{itemize}
    \item гипогликемия;
    \item гипоальбуминемия;
    \item асцит;
    \item нарушение сосудистого тонуса и распределения жидкости;
    \item чувствительность к избытку натрия и объёма;
    \item нарушение метаболизма лактата и других субстратов.
\end{itemize}

\subsection{Практические акценты}
\begin{itemize}
    \item следить за глюкозой чаще;
    \item внимательно оценивать белки и риск третьего пространства;
    \item осторожно относиться к нарастающему положительному балансу;
    \item не забывать о калии;
    \item при асците оценивать не только ``сколько жидкости в животе'', но и реальную перфузию.
\end{itemize}

\subsection{Ловушка}
Асцит и периферические отёки могут создать ложное ощущение ``перегидратации'', хотя эффективный артериальный объём всё ещё может быть недостаточным.

\section{Гипоадренокортицизм: расширенный клинический блок}
\subsection{Почему это важно}
Аддисон --- один из самых клинически ``электролитных'' диагнозов. Его легко пропустить, если смотреть только на отдельные цифры без паттерна.

\subsection{Классический профиль}
\begin{itemize}
    \item гиповолемия;
    \item гипонатриемия;
    \item гиперкалиемия;
    \item слабость;
    \item рвота и/или диарея;
    \item иногда гипогликемия;
    \item нередко преренальная азотемия.
\end{itemize}

\subsection{Важная оговорка}
Электролиты могут быть не ``классическими'' в каждый момент времени. Атипичный гипоадренокортицизм возможен, и даже у таких пациентов минералокортикоидный компонент может развиться позже.

\subsection{Практические выводы}
\begin{itemize}
    \item подозрение на Аддисон не снимается только потому, что калий сейчас ``не очень высокий'';
    \item повторный контроль электролитов может быть нужен даже при неполной картине;
    \item сначала стабилизируют пациента, затем продолжают уточняющую диагностику.
\end{itemize}

\section{Экстренные клинические сценарии}
\subsection{Кот с уретральной обструкцией и гиперкалиемией}
\textbf{Что обычно видим:}
брадикардия, слабость, отсутствие мочеиспускания, переполненный мочевой пузырь, депрессия, гиперкалиемия, азотемия.

\textbf{Что проверить сразу:}
\begin{itemize}
    \item ЭКГ;
    \item \ka, \na, \cl, креатинин;
    \item кислотно-основное состояние;
    \item объёмный статус;
    \item наличие разрыва мочевых путей при сомнительной картине.
\end{itemize}

\textbf{Что делать сначала:}
\begin{enumerate}
    \item стабилизировать миокард кальцием при ЭКГ-изменениях;
    \item инсулин + декстроза для внутриклеточного сдвига \ka;
    \item начать инфузионную коррекцию по состоянию;
    \item быстро устранить обструкцию или декомпрессировать мочевой пузырь по показаниям;
    \item далее --- пересмотреть \ka\ и ЭКГ в динамике.
\end{enumerate}

\textbf{Чего не делать:}
\begin{itemize}
    \item не ждать ``пока сдадутся все анализы'', если есть ЭКГ-угроза;
    \item не считать, что кальций решает проблему гиперкалиемии окончательно;
    \item не забывать, что сбалансированные изотонические растворы допустимы и нередко полезны.
\end{itemize}

\subsection{Собака с гипоадренокортицизмом}
\textbf{Типичная картина:}
слабость, рвота, диарея, гиповолемия, гипонатриемия, гиперкалиемия, иногда гипогликемия.

\textbf{Первые шаги:}
\begin{enumerate}
    \item стабилизация перфузии;
    \item коррекция гиперкалиемии по тяжести;
    \item забор крови на гормональную диагностику, если это не задерживает помощь;
    \item глюкокортикоидная поддержка по протоколу клиники;
    \item дальнейшая коррекция Na и общего водного дефицита после стабилизации.
\end{enumerate}

\textbf{Ловушка:}
не лечить только ``цифры электролитов'', игнорируя общий криз и циркуляторную нестабильность.

\subsection{Пациент с рвотой, гипохлоремией и метаболическим алкалозом}
\textbf{О чём думать:}
\begin{itemize}
    \item потеря HCl;
    \item обструкция верхних отделов ЖКТ;
    \item ``chloride-responsive'' метаболический алкалоз.
\end{itemize}

\textbf{Практика:}
\begin{itemize}
    \item искать источник рвоты/обструкции;
    \item оценивать \cl, \hco, \na, \ka;
    \item рассматривать \(0{,}9\%\) NaCl как один из логичных вариантов коррекции в нужном клиническом контексте.
\end{itemize}

\subsection{ДКА}
\textbf{Ключевая мысль:}
при ДКА дефицит воды сочетается с дефицитом калия, часто фосфора, нередко магния, и всё это может быстро меняться после начала инсулина.

\textbf{Что особенно важно:}
\begin{itemize}
    \item мониторировать \ka\ до, во время и после начала инсулина;
    \item проверять \phos, особенно при слабости или гемолизе;
    \item учитывать натрий с поправкой на гипергликемию;
    \item избегать инерционного ``одинакового'' плана на сутки без частых пересмотров.
\end{itemize}

\subsection{ХБП/ОПП, олигурия и риск перегрузки}
\textbf{Что важно:}
\begin{itemize}
    \item олигурия на фоне продолжающейся инфузии --- опасная зона;
    \item положительный баланс может быстро стать вредным;
    \item гиперкалиемия и гиперфосфатемия требуют более агрессивного мониторинга;
    \item нужно активно исключать обструкцию и пересматривать цели инфузии.
\end{itemize}

\subsection{Эклампсия / симптомная гипокальциемия}
\textbf{Клиника:}
\begin{itemize}
    \item тремор;
    \item тетания;
    \item беспокойство;
    \item гипертермия;
    \item судороги.
\end{itemize}

\textbf{Что делать:}
\begin{enumerate}
    \item медленно вводить кальция глюконат под ЭКГ-контролем;
    \item контролировать дыхание и температуру;
    \item после стабилизации решать вопрос дальнейшей кальциевой поддержки и первичной причины.
\end{enumerate}

\section{Диагностические ловушки}
\begin{enumerate}
    \item \textbf{Псевдогиперкалиемия.} Гемолиз пробы может дать ложное повышение \ka.
    \item \textbf{Натрий при гипергликемии.} Без поправки можно переоценить истинную гипонатриемию.
    \item \textbf{Общий кальций против ионизированного.} При гипоальбуминемии total Ca может вводить в заблуждение.
    \item \textbf{Хлор без поправки на натрий.} Иногда apparent hyper-/hypochloremia лишь отражает сдвиг Na.
    \item \textbf{Нормальный диурез не всегда означает безопасность.} Он может быть на фоне уже нарастающей перегрузки или осмотического диуреза.
    \item \textbf{Ошибка ``лечить цифру, а не гемодинамику''.} Особенно опасна при диснатриемиях.
    \item \textbf{``Нормальная'' ЭКГ не исключает серьёзную гиперкалиемию.} Особенно при быстро меняющемся K.
    \item \textbf{Рефрактерная гипокалиемия без проверки Mg.} Очень частая практическая ошибка.
\end{enumerate}

\section{Роль стационарной команды и среднего персонала}
\subsection{Почему это нужно прописать}
Качество терапии зависит не только от назначений врача, но и от того, насколько чётко команда видит триггеры ухудшения.

\subsection{Что должно быть закреплено в команде}
\textbf{Средний персонал / стационар:}
\begin{itemize}
    \item считает входы и выходы;
    \item взвешивает пациента;
    \item отмечает диурез;
    \item замечает тахипноэ, хемоз, новые отёки, слабость, тремор;
    \item вовремя инициирует повторную ЭКГ или контрольные анализы по триггерам;
    \item сообщает о любом ухудшении раньше, чем оно станет тяжёлым.
\end{itemize}

\textbf{Врач:}
\begin{itemize}
    \item формулирует цель инфузии;
    \item задаёт интервал пересмотра;
    \item прописывает параметры тревоги;
    \item регулярно обновляет план.
\end{itemize}

\subsection{Полезные триггеры для немедленного доклада врачу}
\begin{itemize}
    \item прирост массы тела;
    \item снижение или исчезновение мочи;
    \item новое тахипноэ;
    \item слабость, тремор, судороги;
    \item брадикардия или аритмия;
    \item рвота/диарея, изменившие ожидаемый баланс;
    \item ухудшение сознания.
\end{itemize}

\section{Рабочие алгоритмы}
\subsection{Алгоритм у пациента с подозрением на тяжёлый водно-электролитный сдвиг}
\begin{enumerate}
    \item Оценить перфузию и угрозу жизни.
    \item При гиповолемии --- начать изотоническую ресусцитацию.
    \item Взять минимум лабораторных данных:
    \begin{itemize}
        \item электролиты;
        \item глюкозу;
        \item креатинин/мочевину;
        \item КЩС;
        \item \ica, при необходимости \img\ и \phos.
    \end{itemize}
    \item Оценить диурез и вероятность обструкции.
    \item Определить, какая проблема доминирует:
    \begin{itemize}
        \item объём;
        \item вода;
        \item \na;
        \item \ka;
        \item \cl;
        \item \ca;
        \item \mg;
        \item \phos.
    \end{itemize}
    \item Составить план на ближайшие 4--6 часов и пересматривать его по данным мониторинга.
\end{enumerate}

\subsection{Алгоритм при гипернатриемии}
\begin{enumerate}
    \item Исключить лабораторную ошибку и учесть контекст.
    \item Оценить объёмный статус.
    \item Если есть гиповолемия --- сначала изотоническая стабилизация.
    \item Затем --- медленная коррекция дефицита свободной воды.
    \item Контролировать \na\ серийно.
\end{enumerate}

\subsection{Алгоритм при гиперкалиемии}
\begin{enumerate}
    \item ЭКГ.
    \item Кальций при угрозе аритмии.
    \item Инсулин + декстроза для внутриклеточного сдвига.
    \item Инфузия и восстановление перфузии.
    \item Устранение обструкции/причины.
    \item Диализ при рефрактерном течении.
\end{enumerate}

\subsection{Алгоритм при гипокальциемии}
\begin{enumerate}
    \item Если есть тетания, судороги, выраженный тремор --- не ждать.
    \item Медленно ввести кальция глюконат под ЭКГ-контролем.
    \item Проверить \ica.
    \item Искать причину: эклампсия, гипопаратиреоз, панкреатит, сепсис, этиленгликоль, цитратная нагрузка.
\end{enumerate}

\subsection{Алгоритм при рефрактерной гипокалиемии}
\begin{enumerate}
    \item перепроверить анализ и исключить артефакт;
    \item оценить продолжающиеся потери;
    \item проверить \mg;
    \item пересмотреть инфузионный состав и энтеральное поступление;
    \item проверить кислотно-основное состояние;
    \item исключить избыточную минералокортикоидную активность или продолжающийся осмотический диурез.
\end{enumerate}

\section{Таймлайн пересмотра пациента: 1 час / 4 часа / 12 часов}
\subsection{Через 1 час}
\begin{itemize}
    \item ответил ли пациент на болюс / старт инфузии;
    \item изменились ли CRT, пульс, ЧСС, ментальный статус;
    \item нет ли ухудшения дыхания;
    \item есть ли моча;
    \item нужна ли повторная ЭКГ или срочная коррекция электролита.
\end{itemize}

\subsection{Через 4 часа}
\begin{itemize}
    \item соответствует ли текущий раствор новой клинической задаче;
    \item пересмотрены ли ongoing losses;
    \item изменились ли \na, \ka, \cl, глюкоза, КЩС;
    \item нет ли признаков перегрузки;
    \item не пора ли уменьшать или менять инфузию.
\end{itemize}

\subsection{Через 12 часов}
\begin{itemize}
    \item куда пациент двигается в целом;
    \item сохраняется ли исходная проблема или уже доминирует новая;
    \item нужен ли другой уровень мониторинга;
    \item возможно ли перейти на более безопасную и простую схему;
    \item всё ли ещё требуется из интенсивной коррекции.
\end{itemize}

\important{Повторная оценка через несколько часов почти всегда полезнее, чем самый красивый стартовый план.}

\section{Частые ошибки}
\begin{enumerate}
    \item Корректировать натрий, не восстановив перфузию.
    \item Смотреть только на клиническую ``сухость'', игнорируя массу тела и диурез.
    \item Недооценивать хлор.
    \item Лечить гипокалиемию, не проверив магний.
    \item Оценивать кальций только по общему кальцию при гипоальбуминемии.
    \item Забывать про фосфор при ДКА, рефидинге и тяжёлой анорексии.
    \item Продолжать инфузию на фоне нарастающего положительного баланса без пользы.
\end{enumerate}

\section{Короткая памятка}
\begin{longtable}{>{\raggedright\arraybackslash}p{0.33\textwidth} >{\raggedright\arraybackslash}p{0.62\textwidth}}
\toprule
\textbf{Ситуация} & \textbf{Что важно помнить} \\
\midrule
Гиповолемия & Сначала перфузия, потом тонкая коррекция электролитов \\
Гипернатриемия & Медленная коррекция, риск отёка мозга при слишком быстром снижении \\
Гипонатриемия & Медленная коррекция хронических случаев, риск осмотического повреждения при слишком быстром повышении \\
Гиперкалиемия & ЭКГ, кальций, инсулин с декстрозой, устранить причину \\
Гипокалиемия & Проверить магний, не превышать безопасную скорость в/в коррекции \\
Гипохлоремия & Часто связана с рвотой и метаболическим алкалозом \\
Гиперкальциемия & Ищите неоплазию, гиперпаратиреоз, гипоадренокортицизм, ХБП \\
Гипокальциемия & Для решения нужен ионизированный кальций, а не только общий \\
Гипофосфатемия & Помнить про ДКА, рефидинг, гемолиз, мышечную слабость \\
Перегрузка жидкостью & Положительный баланс без пользы --- не повод увеличивать капельницу \\
\bottomrule
\end{longtable}

\section{Приложение 1: практические численные ориентиры}
\begin{itemize}
    \item Диурез: около \(1\ \text{мл/кг/ч}\) и выше как практический ориентир.
    \item Максимальная скорость изменения натрия: примерно \(0{,}5\ \text{ммоль/л/ч}\).
    \item Часто используемый суточный предел изменения натрия при хронических нарушениях: около \(8\text{--}12\ \text{ммоль/л/сут}\).
    \item Верхний ориентир скорости внутривенного введения калия: около \(0{,}5\ \text{ммоль/кг/ч}\).
\end{itemize}

\section{Приложение 2: лист водного баланса на 24 часа}
\renewcommand{\arraystretch}{1.25}
\begin{longtable}{|p{0.08\textwidth}|p{0.18\textwidth}|p{0.18\textwidth}|p{0.18\textwidth}|p{0.18\textwidth}|p{0.12\textwidth}|}
\hline
\textbf{Час} & \textbf{Инфузия / болюсы} & \textbf{Перорально / корм} & \textbf{Моча} & \textbf{Прочие потери} & \textbf{Баланс} \\
\hline
00--01 & & & & & \\
\hline
01--02 & & & & & \\
\hline
02--03 & & & & & \\
\hline
03--04 & & & & & \\
\hline
04--05 & & & & & \\
\hline
05--06 & & & & & \\
\hline
06--07 & & & & & \\
\hline
07--08 & & & & & \\
\hline
08--09 & & & & & \\
\hline
09--10 & & & & & \\
\hline
10--11 & & & & & \\
\hline
11--12 & & & & & \\
\hline
12--13 & & & & & \\
\hline
13--14 & & & & & \\
\hline
14--15 & & & & & \\
\hline
15--16 & & & & & \\
\hline
16--17 & & & & & \\
\hline
17--18 & & & & & \\
\hline
18--19 & & & & & \\
\hline
19--20 & & & & & \\
\hline
20--21 & & & & & \\
\hline
21--22 & & & & & \\
\hline
22--23 & & & & & \\
\hline
23--24 & & & & & \\
\hline
\end{longtable}
\renewcommand{\arraystretch}{1}

\section{Приложение 3: чек-лист пересмотра инфузии}
\begin{itemize}
    \item Улучшилась ли перфузия?
    \item Что происходит с массой тела?
    \item Какой почасовой диурез?
    \item Есть ли новые отёки, хемоз, тахипноэ, выпоты?
    \item Что произошло с \na, \ka, \cl, \ica, \img, \phos?
    \item Изменилось ли КЩС?
    \item Текущий раствор всё ещё соответствует цели?
    \item Не лечу ли я уже ``старую'' проблему, которая изменилась несколько часов назад?
\end{itemize}

\section{Приложение 4: одностраничная памятка для ОРИТ}
\begin{enumerate}
    \item Сначала перфузия, потом тонкая коррекция.
    \item Не корректировать натрий быстро.
    \item Гиперкалиемия = ЭКГ + кальций + внутриклеточный сдвиг + устранение причины.
    \item Рефрактерная гипокалиемия --- проверь магний.
    \item При подозрении на клинически значимую гипокальциемию ориентируйся на ионизированный кальций.
    \item При ДКА и рефидинге не забывай про фосфор.
    \item Положительный баланс без пользы --- пересмотри инфузию.
    \item Нормальный общий кальций не исключает проблему \ica.
    \item Хлор интерпретируй вместе с натрием и КЩС.
    \item Повторная оценка через несколько часов важнее красивого плана ``на сутки''.
\end{enumerate}

\section*{Источники для врача}
\begin{enumerate}
    \item AAHA Fluid Therapy Guidelines for Dogs and Cats, 2024.
    \item MSD / Merck Veterinary Manual: fluid therapy, fluid overload, potassium disorders, calcium disorders, eclampsia, monitoring of the critically ill small animal.
    \item Современные обзоры и исследования по хлору, магнию, диснатриемиям и электролитным нарушениям у собак и кошек.
\end{enumerate}

\vspace{1em}
\noindent\textcolor{textgray}{\small
Замечание: численные пороги и интервалы мониторинга следует сверять с референсами вашей лаборатории, используемым анализатором и клиническим контекстом пациента. Данный документ предназначен как рабочее руководство, а не как замена клиническому суждению.}

\end{document}